\documentclass[12pt]{article}
\usepackage[T1]{fontenc}
\usepackage[utf8]{inputenc}
\usepackage{lmodern}
\usepackage{textcomp}
\usepackage{enumitem}
\usepackage{geometry}
\geometry{margin=1in}
\begin{document}
\begin{center}
Answer Key - Economics - K-12 Level Assessment 20 \
\small (Answers are ordered exactly as in the question paper.)
\end{center}

\section*{Multiple Choice}
\begin{enumerate}[label=\arabic*.]
\item \textbf{Answer:} B
\\ \textit{Options:} A) Net exports fall.; B) Income is transferred from domestic steel consumers to domestic steel producers.; C) Allocative efficiency is improved.; D) Income is transferred from domestic steel to foreign steel producers.
\\ \textit{Note:} A protective tariff on imported steel raises the price of foreign steel, leading consumers to pay more for steel products while benefiting domestic producers by reducing foreign competition and allowing them to charge higher prices.
\item \textbf{Answer:} D
\\ \textit{Options:} A) Both supply and demand shift rightward.; B) Both supply and demand shift leftward.; C) Supply shifts to the right; demand shifts to the left.; D) Supply shifts to the left; demand shifts to the right.
\\ \textit{Note:} The correct answer is based on the fundamental principle of supply and demand, where a leftward shift in supply indicates a decrease in quantity supplied, while a rightward shift in demand indicates an increase in quantity demanded, both of which together lead to a higher equilibrium price for corn.
\item \textbf{Answer:} B
\\ \textit{Options:} A) the value of the dollar will tend to appreciate.; B) the value of the dollar will tend to depreciate.; C) exchange rates will be affected but not the value of the dollar.; D) the exchange rate will not be affected.
\\ \textit{Note:} When prices rise in the United States relative to other countries, it makes U.S. goods more expensive for foreign buyers, leading to a decrease in demand for the dollar and causing it to depreciate in value.
\item \textbf{Answer:} C
\\ \textit{Options:} A) Demand is downward-sloping.; B) The demand curve lies above the marginal revenue curve.; C) Price is determined by the equilibrium in the entire market.; D) Average revenue differs from price.
\\ \textit{Note:} In a perfectly competitive market, individual firms are price takers, meaning the market price is established by the overall supply and demand dynamics in the market, rather than by any single firm's decisions.
\item \textbf{Answer:} A
\\ \textit{Options:} A) There would be no change in the amount of corn demanded or supplied.; B) There would be a shortage created of corn.; C) There would be a surplus created of corn.; D) The producers of corn would lose revenue due to the decreased price.
\\ \textit{Note:} A price ceiling set above the equilibrium price does not affect the market because it does not restrict the price from reaching its natural level, resulting in no change in the quantity of corn demanded or supplied.
\end{enumerate}

\section*{True / False}
\begin{enumerate}[label=\arabic*.]
\setcounter{enumi}{5}
\item \textbf{Answer:} False
\\ \textit{Options:} A) True; B) False
\\ \textit{Note:} The statement is false because when demand exceeds supply, suppliers are generally incentivized to increase production or prices rather than automatically decreasing the quantity supplied.
\item \textbf{Answer:} True
\\ \textit{Options:} A) True; B) False
\\ \textit{Note:} An increase in the standard of living signifies that the average income or wealth per person has risen, which occurs when economic output grows at a faster rate than the population, allowing individuals to enjoy more goods and services.
\item \textbf{Answer:} False
\\ \textit{Options:} A) True; B) False
\\ \textit{Note:} An increasing national debt can lead to concerns about the country's financial stability, which can decrease demand for U.S. dollars and cause the dollar's value to depreciate rather than appreciate.
\item \textbf{Answer:} False
\\ \textit{Options:} A) True; B) False
\\ \textit{Note:} Rising prices can erode purchasing power, meaning that even if nominal household incomes increase, consumers may not be able to afford more goods and services if their real income remains stagnant or decreases.
\item \textbf{Answer:} False
\\ \textit{Options:} A) True; B) False
\\ \textit{Note:} National income is not always greater than GDP because GDP measures the total value of all goods and services produced within a country, while national income accounts for income earned by residents, which can be impacted by factors like depreciation and net income from abroad.
\end{enumerate}

\section*{Long Form Response}
\begin{enumerate}[label=\arabic*.]
\setcounter{enumi}{10}
\item \textbf{Answer:} The cost of a Big Mac varies significantly across countries due to factors such as exchange rates, local purchasing power, and economic conditions. In this analysis, Mexico has the highest cost for a Big Mac at \$5. This elevated price can be attributed to the combination of local inflation, the cost of ingredients, and the effects of currency exchange rates that make imported goods relatively more expensive. Compared to the United States, England, and China, where prices are often lower, Mexico's economic situation leads to a higher retail price for this popular fast-food item. Understanding these differences helps illustrate how global economics influences everyday purchases.
\\ \textit{Note:} correct because it identifies that the high cost of a Big Mac in Mexico is influenced by local inflation, ingredient costs, and unfavorable exchange rates, which collectively elevate the price compared to the other countries analyzed.
\item \textbf{Answer:} The Federal Reserve can use bond purchases as a policy tool to address an economic recession by increasing the money supply in the economy. When the Fed buys bonds, it pays for them with newly created money, which banks receive. This action lowers interest rates, making borrowing cheaper for consumers and businesses. As a result, lower interest rates encourage spending and investment, which can help stimulate economic growth. Overall, bond purchases can lead to increased economic activity, helping to pull the economy out of a recession.
\\ \textit{Note:} correct because it accurately describes how bond purchases by the Federal Reserve increase the money supply, lower interest rates, and stimulate spending and investment, thereby helping to counteract an economic recession.
\end{enumerate}

\vfill
\noindent\textit{End of Answer Key}
\end{document}